\section{Experiments}
\label{sec:experiments}

\subsection{Protocol}

We construct an illustrative benchmark from ImageNet-C~\citep{hendrycks2019robustness},
Cityscapes-C, and ACDC~\citep{sakaridis2021acdc}.  A ViT-B/16 backbone is
adapted online with batch size 16.  We set $B=0.45$, adapter ratio
$\gamma=0.25$, temperature $\tau=0.7$, and update with AdamW at
$2\times10^{-5}$.  Baselines are source-only inference, TENT, EATA, and full
adaptation using the same unsupervised loss.  Classification uses top-1
accuracy; segmentation uses mean intersection-over-union (mIoU).  Latency is
measured after warm-up on one illustrative accelerator.  All values in this
section are synthetic and serve only to demonstrate reporting conventions.

\begin{table}[t]
  \centering
  \caption{Illustrative accuracy and efficiency.  Higher accuracy/mIoU is
  better; lower latency is better.  Best values are bold.}
  \label{tab:main}
  \resizebox{\columnwidth}{!}{%
  \begin{tabular}{lcccc}
    \toprule
    Method & IN-C $\uparrow$ & ACDC $\uparrow$ & Tokens $\downarrow$ & ms $\downarrow$\\
    \midrule
    Source only & 54.8 & 52.1 & 100\% & 18.4\\
    TENT~\citep{wang2021tent} & 60.7 & 55.6 & 100\% & 27.9\\
    EATA~\citep{niu2022efficient} & 62.3 & 57.1 & 100\% & 25.8\\
    Full adaptation & \textbf{64.1} & \textbf{59.0} & 100\% & 29.6\\
    \rowcolor{palegreen}
    \method{} & 63.3 & 58.4 & \textbf{43\%} & \textbf{19.5}\\
    \bottomrule
  \end{tabular}}
\end{table}

\subsection{Main results}

\Cref{tab:main} shows the central trade-off.  \method{} recovers 8.5 points
over source-only inference on ImageNet-C and approaches full adaptation within
0.8 points, while updating fewer than half of all tokens.  The latency gain is
smaller than the token reduction because patch encoding, routing, and the task
head remain dense.  On ACDC, the 0.6-point gap to full adaptation suggests that
the router still identifies regions affected by adverse weather.

\begin{figure}[t]
  \centering
  \begin{subfigure}[b]{0.48\columnwidth}
    \centering
    \begin{tikzpicture}[x=0.36cm,y=0.36cm]
      \fill[mist] (0,0) rectangle (8,5);
      \fill[ink!75] (0,0) -- (8,0) -- (6.2,2.1) -- (1.8,2.1) -- cycle;
      \fill[cvprblue!35] (0,2.1) rectangle (8,5);
      \fill[white!85] (0.5,3.9) circle (0.35);
      \foreach \x/\y/\c in {0.5/0.5/accent,1.5/0.5/accent,2.5/0.5/gold,
      3.5/0.5/coral,4.5/0.5/coral,5.5/0.5/gold,6.5/0.5/accent,7.5/0.5/accent,
      0.5/1.5/accent,1.5/1.5/gold,2.5/1.5/coral,3.5/1.5/coral,
      4.5/1.5/coral,5.5/1.5/coral,6.5/1.5/gold,7.5/1.5/accent,
      0.5/2.5/accent,1.5/2.5/accent,2.5/2.5/gold,3.5/2.5/gold,
      4.5/2.5/gold,5.5/2.5/gold,6.5/2.5/accent,7.5/2.5/accent}
      {\fill[\c,opacity=.42] (\x-.43,\y-.43) rectangle (\x+.43,\y+.43);}
    \end{tikzpicture}
    \caption{Localized road shift}
  \end{subfigure}\hfill
  \begin{subfigure}[b]{0.48\columnwidth}
    \centering
    \begin{tikzpicture}[x=0.36cm,y=0.36cm]
      \fill[cvprblue!20] (0,0) rectangle (8,5);
      \fill[ink!70] (0,0) rectangle (8,1.8);
      \foreach \x in {0.5,1.5,...,7.5}
        \foreach \y in {0.5,1.5,...,4.5}
          {\pgfmathparse{int(mod(2*\x+3*\y,5))}
           \ifnum\pgfmathresult<2
             \fill[coral,opacity=.40] (\x-.43,\y-.43) rectangle (\x+.43,\y+.43);
           \else
             \fill[gold,opacity=.32] (\x-.43,\y-.43) rectangle (\x+.43,\y+.43);
           \fi}
    \end{tikzpicture}
    \caption{Global sensor noise}
  \end{subfigure}
  \vspace{-2pt}
  \begin{minipage}{\columnwidth}
    \centering\scriptsize
    \textcolor{accent}{\rule{7pt}{7pt}} reuse \quad
    \textcolor{gold}{\rule{7pt}{7pt}} adapter \quad
    \textcolor{coral}{\rule{7pt}{7pt}} full block
  \end{minipage}
  \caption{Illustrative routing maps.  Compute follows a localized shift
  (left), but widespread noise pushes many tokens to expensive paths (right).}
  \label{fig:routing}
\end{figure}

\subsection{Ablation and budget sweep}

Removing augmentation agreement increases confident mistakes; removing
novelty causes the cache to persist after abrupt scene changes
(\cref{tab:ablation}).  A global threshold saves compute but under-serves rare
classes.  The complete router gives the strongest accuracy at similar cost.

\begin{table}[t]
  \centering
  \caption{Synthetic ablation on ImageNet-C.  ``Full route'' is the fraction
  of tokens assigned to the most expensive path.}
  \label{tab:ablation}
  \begin{tabular}{lccc}
    \toprule
    Variant & Acc. $\uparrow$ & ECE $\downarrow$ & Full route $\downarrow$\\
    \midrule
    Global threshold & 61.8 & 6.4 & 24\%\\
    $-$ agreement & 62.1 & 7.0 & 22\%\\
    $-$ novelty & 62.5 & 5.9 & 20\%\\
    \rowcolor{palegreen}
    Complete \method{} & \textbf{63.3} & \textbf{4.7} & \textbf{19\%}\\
    \bottomrule
  \end{tabular}
\end{table}

Varying $B$ from 0.30 to 0.70 yields accuracies of 61.9, 63.3, and 63.8,
respectively.  The diminishing return above $B=0.45$ indicates that the router
ranks useful tokens rather than uniformly degrading the network.

% -----------------------------------------------------------------------------
